% !TEX root = task-3-4-2.tex
% Editable task bank. Append new BankTask blocks at the end of the document.
\documentclass[12pt,a4paper]{article}
\usepackage{physics-textbook}
\usepackage{graphicx,needspace}
\usetikzlibrary{decorations.pathmorphing}
\newif\ifShowSolutions
\ifdefined\TutorEdition
  \ShowSolutionstrue
\else
  \ShowSolutionsfalse % Change to \ShowSolutionstrue to display solutions.
\fi
\newcounter{banktask}
\newsavebox{\BankTaskBox}
\newcommand{\BankTaskContent}[3]{%
  \noindent\textbf{\thebanktask.}\quad #2\par
  \ifShowSolutions
    \Needspace{50mm}% Keep the solution heading with its opening text or diagram.
    #3\par
    \ifdefined\TutorEdition
      \ifcsname TutorSource#1\endcsname
        \nopagebreak\smallskip{\small\csname TutorSource#1\endcsname\par}%
      \fi
    \fi
  \fi
}
\newcommand{\BankTask}[3]{%
  \refstepcounter{banktask}\par\medskip
  \sbox{\BankTaskBox}{\begin{minipage}{\linewidth}%
    \BankTaskContent{#1}{#2}{#3}\end{minipage}}%
  \ifdim\dimexpr\ht\BankTaskBox+\dp\BankTaskBox\relax>\dimexpr\textheight-20mm\relax
    \Needspace{35mm}\BankTaskContent{#1}{#2}{#3}%
  \else
    \Needspace{\dimexpr\ht\BankTaskBox+\dp\BankTaskBox\relax}%
    \noindent\usebox{\BankTaskBox}\par
  \fi
  \medskip
}
\begin{document}
\BookSubsection{3.4.2}{Потенциальность электростатического поля}


% Source E5F230: https://ege.fipi.ru/bank/questions.php?search=1&proj=BA1F39653304A5B041B656915DC36B38&qid=E5F230
\BankTask{E5F230}{
Однородное электростатическое поле создано равномерно заряженной протяженной горизонтальной пластиной. Линии напряжённости поля направлены вертикально вверх (см. рисунок).
\begin{enumerate}
\item Пластина имеет отрицательный заряд.
\item Потенциал электростатического поля в точке $A$ ниже, чем в точке $C$.
\item Работа электростатического поля по перемещению пробного точечного отрицательного заряда из точки $A$ в точку $B$ отрицательна.
\item Если в точку $B$ поместить пробный точечный положительный заряд, то на него со стороны пластины будет действовать сила, направленная вертикально вверх.
\item Напряжённость поля в точках $A$, $B$ и $C$ одинакова.
\end{enumerate}
\begin{center}
\begin{tikzpicture}[book diagram]
\draw[line width=1.2pt] (-0.45,0)--(2.05,0);
\foreach \x in {0,0.55,1.1,1.65}{
\draw (\x,0)--(\x,2.55);
\draw[-{book tip}] (\x,1.65)--(\x,2.03);
}
\draw[dashed] (0,0.8)--(1.65,0.8);
\draw[dashed] (0,0.8)--(1.65,1.68);
\draw (1.50,0.8)--(1.50,0.95)--(1.65,0.95);
\fill (0,0.8) circle (2pt) node[left=3pt] {$A$};
\fill (1.65,0.8) circle (2pt) node[right=3pt] {$B$};
\fill (1.65,1.68) circle (2pt) node[right=3pt] {$C$};
\node at (1.35,2.48) {$\vec E$};
\end{tikzpicture}
\end{center}
}{
\smallskip\noindent\textbf{Решение от GPT}\par
1) Неверно: линии поля направлены от пластины, следовательно, её заряд положителен. 2) Неверно: потенциал убывает вдоль поля, поэтому $\varphi_A>\varphi_C$. 3) Неверно: точки $A$ и $B$ лежат на одной эквипотенциальной поверхности, так что $A_{AB}=q(\varphi_A-\varphi_B)=0$. 4) Верно: для положительного заряда сила $\vec F=q\vec E$ направлена вверх. 5) Верно: в однородном поле $\vec E_A=\vec E_B=\vec E_C$. Ответ: $4,\ 5$.
}

\end{document}
